\subsection{上积分}

定义 11. --- 对每个函数 \(f \in \mathcal{F}_{+}(T)\) ，定义 f 的上积分（关于测度 \(\mu\) ）为有限或无限的正数

\[
\mu^ {*} (f) = \inf _ {g} \mu^ {\bullet} (g),\tag{9}
\]

其中 \(g\) 取遍 \(\geqslant f\) 的下半连续函数全体。也使用记号 \(\int^{*} f(t) d\mu(t)\) 与 \(\int^{*} f d\mu\) 。当 \(T\) 局部紧时，这一定义与通常的定义一致 (Ch. V, §1, No.~1, Prop. 4)。显然 \(\mu^{\bullet}(f) \leqslant \mu^{*}(f)\) ，且当 \(f\) 下半连续时取等号。若 \(A\) 是 \(T\) 的子集，则用 \(\mu^{*}(A)\) 代替 \(\mu^{*}(\varphi_{A})\) ，这个数称为 \(A\) 的外测度。具有有限外测度的可测集称为可积集，与局部紧空间的情形一样。

函数 \(\mathbf{f}\) 若以 Banach 空间或 \(\overline{\mathbf{R}}\) 为值且满足 \(\mu^{*}(|\mathbf{f}|)=0\) ，就称为可忽略的；集 \(A\subset T\) 称为可忽略的，如果 \(\varphi_{A}\) 可忽略，即 \(\mu^{*}(A)=0\) 。几乎处处这一说法按 Ch. IV, §2, No.~3 的方式引入。

命题 12. --- 函数 \(\mu^{*}\) 是 T 上的一个负担。

No.~1 的 Def. 1 的性质 \(a\) 、\(b\) 、\(c\) 是明显的。性质 \(d\) 的证明与 Ch. IV, §1, No.~3 的 Th. 3 的证明相同，只需顾及 Props. 4 与 5 \(a\) )。

推论. --- 以 Banach 空间或 \(\overline{R}\) 为值的函数 f 可忽略，当且仅当 \(\mathbf{f}(t)=0\) 几乎处处成立。

立即化归为正函数的情形。此后的证明与 Ch. IV, §2, No.~3 的 Th. 1 的证明相同。

命题 13. --- 对 T 的每个子集 A，\(\mu^{*}(\mathbf{A})\) 是包含 A 的诸开集的外测度的下确界。

证明与 Ch. IV, §1, No.~4 的 Prop. 19 的证明相同。

定义 12. --- 设 f 是定义在 T 上、以 Banach 空间或 \(\overline{R}\) 为值的函数。若 f 在可数个可积开集之并的余集上为零，则称 f 对测度 \(\mu\) 适度，或 \(\mu\) -适度。若函数 \(\varphi_{A}\) 适度，则称 T 的子集 A 适度。若函数 1 是 \(\mu\) -适度的，则称测度 \(\mu\) 适度。

例如，由于负担 \(\mu^{\bullet}\) 局部有界，T 的每个紧子集 K 都含于一个使 \(\mu^{\bullet}(V)<+\infty\) 的开集 V；因此在紧集外为零的函数是适度的。可忽略函数是适度的。Ch. V, §1, No.~2 中 Def. 2 后面的注记可立即推广到当前情形。特别地，一列适度正函数的和是适度的。

注记. --- 1) 在 Lindelöf 空间 T 上 (TG, IX, Appendix I, Def. 1)，\(^{1}\) 特别地在 Souslin 空间上 (ibid., Prop. 1 的 Cor.)，每个测度都是适度的。因为有限测度的开集构成 T 的一个覆盖，从中可以取出 T 的一个可数覆盖。

\begin{enumerate}
\def\labelenumi{\arabic{enumi})}
\setcounter{enumi}{1}
\tightlist
\item
  然而要注意，存在一列对 \(\mu\) 具有有限测度、并以 T 为并的 Borel 集，并不一定蕴含存在一列以 T 为并、具有限测度的开集（换句话说，不蕴含 \(\mu\) 适度）。见 Exer. 8。
\end{enumerate}

命题 14. --- 设 \(f \in \mathcal{F}_{+}(\mathrm{T})\) 。若 \(f\) 是 \(\mu\) -适度的，则 \(\mu^{*}(f) = \mu^{\bullet}(f)\) ；若 \(f\) 不是 \(\mu\) -适度的，则 \(\mu^{*}(f) = +\infty\) 。

若 \(\mu^{*}(f) < +\infty\) ，则存在下半连续函数 \(g \geqslant f\) 使 \(\mu^{\bullet}(g) < +\infty\) 。对每个 \(n \in \mathbf{N}\) ，设 \(\mathbf{G}_{n}\) 是 \(t \in \mathbf{T}\) 中使 \(g(t) > 1/n\) 者的全体；集 \(\mathbf{G}_{n}\) 是开集，有 \(\mu^{\bullet}(\mathbf{G}_{n}) \leqslant n\mu^{\bullet}(g) < +\infty\) ，且 \(f\) 在诸 \(\mathbf{G}_{n}\) 之并外为零：因此函数 \(f\) 适度。

其次，证明 \(\mu^{*}\) 与 \(\mu^{\bullet}\) 在适度函数上取相同的值。由于 \(\mu^{*}\) 与 \(\mu^{\bullet}\) 都是负担，只需对如下情形建立关系 \(\mu^{*}(f)=\mu^{\bullet}(f)\)：f 是正函数，被常数 M 控制且在一个具有限测度的开集 G 外为零；我们现在就来证明这一点。

测度 \(\mu\) 是 \(\mathcal{M}(\mathrm{T})\) 中一个由具紧支集的测度组成的递增定向族 \((\mu_i)_{i\in I}\) 的上确界：这由 No.~8 的 Prop. 9 立即得出。设 \(g\) 是 \(\mathrm{T}\) 上介于 \(f\) 与下半连续函数 \(\mathrm{M}\varphi_{\mathrm{G}}\) 之间的下半连续函数。令 \(\nu_i = \mu - \mu_i\) ；则 \(\mu^{\bullet} = \mu_i^{\bullet} + \nu_i^{\bullet}\) (No.~2, 注 1)，因此

\[
\begin{array}{r l} \mu^ {\bullet} (g) - \mu^ {\bullet} (f) & = (\mu_ {i} ^ {\bullet} (g) - \mu_ {i} ^ {\bullet} (f)) + (\nu_ {i} ^ {\bullet} (g) - \nu_ {i} ^ {\bullet} (f)) \\ & \leqslant (\mu_ {i} ^ {\bullet} (g) - \mu_ {i} ^ {\bullet} (f)) + \nu_ {i} ^ {\bullet} (\mathrm{M} \varphi_ {\mathrm{G}}). \end{array}
\]

有 \(\nu_{i}^{\bullet}(\mathrm{M}\varphi_{\mathrm{G}}) = \mu^{\bullet}(\mathrm{M}\varphi_{\mathrm{G}}) - \mu_{i}^{\bullet}(\mathrm{M}\varphi_{\mathrm{G}})\) 且 \(\mu^{\bullet}(\mathrm{M}\varphi_{\mathrm{G}}) = \sup \mu_{i}^{\bullet}(\mathrm{M}\varphi_{\mathrm{G}})\) (No.~7, Prop. 6)；因此数 \(\nu_{i}^{\bullet}(\mathrm{M}\varphi_{\mathrm{G}})\) 可通过适当选取 \(i\) 而任意小。于是问题归结为证明：对任意数 \(c > 0\) 与任意指标 \(i \in I\) ，能找到下半连续函数 \(g\) ，介于 \(f\) 与 \(\mathrm{M}\varphi_{\mathrm{G}}\) 之间，使 \(\mu_{i}^{\bullet}(g) - \mu_{i}^{\bullet}(f) \leqslant c\) 。现设 \(L\) 是测度 \(\mu_{i}\) 的紧支集，\(\lambda\) 是测度 \((\mu_{i})_{L}\) ；由于 \(\mu_{i}\) 集中于 \(L\) ，有 \(\mu_{i}^{\bullet}(h) = \mu_{i}^{\bullet}(h\varphi_{L}) = \lambda^{\bullet}(h_{L})\) （对每个函数 \(h \in \mathcal{F}_{+}(T)\)）(No.~1, Lemma 1 与 No.~2, Prop. 2)；因此

\[
\mu_ {i} ^ {\bullet} (g) - \mu_ {i} ^ {\bullet} (f) = \lambda^ {\bullet} (g _ {\mathrm{L}}) - \lambda^ {\bullet} (f _ {\mathrm{L}}).
\]

但 L 是紧的；因此 \(\lambda^{\bullet} = \lambda^{*}\) ，于是存在定义在 L 上的下半连续函数 h，它 \(\geqslant f_{L}\) 且使 \(\lambda^{\bullet}(h) \leqslant \lambda^{\bullet}(f_{\mathrm{L}}) + c\) 。由于集 L 在 T 中闭，在 L 上等于 h 而在 T - L 上等于 \(+\infty\) 的函数 k 在 T 上下半连续且 \(\geqslant f\) ，并且 \(\lambda^{\bullet}(k_{\mathrm{L}}) = \lambda^{\bullet}(h) \leqslant \lambda^{\bullet}(f_{\mathrm{L}}) + c\) 。只需令 \(g = \inf(k, M\varphi_{G})\) ：g 下半连续，\(g \geqslant f\) ，且

\[
\mu_ {i} ^ {\bullet} (g) \leqslant \mu_ {i} ^ {\bullet} (k) = \lambda^ {\bullet} (k _ {\mathrm{L}}) \leqslant \lambda^ {\bullet} (f _ {\mathrm{L}}) + c = \mu_ {i} ^ {\bullet} (f) + c.
\]

推论 1. --- 一个函数可忽略，必须且只需它局部可忽略且适度。

推论 2. --- 函数 f 局部可忽略，必须且只需每个 \(x \in T\) 都有邻域 V 使 \(f\varphi_{V}\) 可忽略。

因为若这一性质满足，则对每个紧集 K，\(f\varphi_{K}\) 可忽略，因此 f 局部可忽略 (No.~2, Prop. 2)。反之，设 f 局部可忽略，x 是 T 的一点；x 有一个具有限测度的开邻域 V。于是函数 \(f\varphi_{V}\) 局部可忽略且适度，从而可忽略。

推论 3. --- 设 \(\mathbf{f}\) 是定义在 \(\mathbf{T}\) 上的适度函数。存在两两不相交的紧集序列 \((\mathbf{K}_n)\) 与可忽略集 \(\mathbf{H}\) ，使 \(\mathbf{f} = \mathbf{f}\varphi_{\mathrm{H}} + \sum_{n} \mathbf{f}\varphi_{\mathrm{K}_n}\) 。

因为设 G 是可数个可积开集之并的集，f 在 G 外为零；则 G 是两两不相交的紧集序列 \((K_{n})\) 与一个局部可忽略集 H 的并 (No.~8, Prop. 11)；而 H 适度，因此可忽略。

推论 4. --- 若 \(\mu\) 与 \(\nu\) 是 T 上两个使 \(\mu^{*} = \nu^{*}\) 的测度，则 \(\mu = \nu\) 。

因为等式 \(\mu^{*}=\nu^{*}\) 蕴含 \(\mu^{\bullet}(f)=\nu^{\bullet}(f)\) 对每个对 \(\mu\) 与 \(\nu\) 适度的正函数 f 成立，从而对每个具紧支集的正函数成立；由此得 \(\mu^{\bullet}=\nu^{\bullet}\) (No.~2, Prop. 2)，再得 \(\mu=\nu\) (No.~2, Prop. 2 的 Cor.)。

推论 5. --- 若 \(\mu\) 是 T 上的适度测度，则存在一列具紧支集的测度 \((\mu_n)_{n \in \mathbf{N}}\) 使 \(\mu = \sum_{n \in \mathbf{N}} \mu_n\) 。

根据假设，常函数 1 是 \(\mu\) -适度的。把 Cor. 3 应用于 \(f = 1\) 的情形；于是存在 T 的两两不相交的紧子集序列 \((\mathrm{K}_n)_{n \in \mathbf{N}}\) ，使 \(1 = \sum_{n \in \mathbf{N}} \varphi_{\mathrm{K}_n}\) \(\mu\) -几乎处处成立。

设 \(\mu_{n}\) 是由测度 \(\mu_{\mathrm{K}_n}\) （在 \(\mathrm{K}_n\) 上）定义的 T 上的测度 (No.~3, 例 2)。已知 (No.~6, 注 2) \(\mu_{n}\) 具紧支集，且 \(\mu_n^\bullet (f) = \mu^\bullet (f\varphi_{\mathrm{K}_n})\) （对 \(f\in \mathcal{F}_{+}(\mathrm{T})\)）。而 \(f\) 等于 \(\sum_{n\in \mathbf{N}}f\varphi_{\mathrm{K}_n}\) （\(\mu\) -几乎处处），由此

\[
\mu^ {\bullet} (f) = \sum_ {n \in \mathbf {N}} \mu^ {\bullet} (f \varphi_ {\mathrm{K} _ {n}}) = \sum_ {n \in \mathbf {N}} \mu_ {n} ^ {\bullet} (f).
\]

于是 \(\mu = \sum_{n\in \mathbf{N}}\mu_n\) (No.~7, Prop. 7)。

